\subsection*{S1a. Numbers: integer versus fractional}

JSON has one number type; OCaml distinguishes \texttt{int} from
\texttt{float}. A path holding numbers must be assigned one of them, and the
same path may hold both integral and fractional values across the corpus.

\begin{lstlisting}[style=json]
{ "qty": 1 }      { "qty": 2 }        -- integral everywhere
{ "price": 1 }    { "price": 2.5 }    -- mixed at one path
\end{lstlisting}

\options
\begin{enumerate}[label=\textbf{(\alph*)}]
  \item \textbf{Split by value, widen on mix.} A path is \texttt{int} if every
        value observed there is integral, \texttt{float} otherwise. Under this
        reading \texttt{1.0} is integral and yields \texttt{int}.
  \item \textbf{Split by token, widen on mix.} Follow the lexical form the
        parser already reports: \texttt{1} is an integer, \texttt{1.0} is a
        float regardless of its value. A path carrying any float becomes
        \texttt{float}. \selected
  \item \textbf{Every number is float.} Map all JSON numbers to \texttt{float},
        matching the IEEE-double reading most JSON implementations take. No
        widening logic, but counts, indices and identifiers all become floats.
  \item \textbf{Split, and reject on mix.} Keep \texttt{int} and \texttt{float}
        separate, and treat a path carrying both as a hard error rather than
        widening it.
\end{enumerate}

\decision{Split by token, widen on mix. A path whose values are all written
without a decimal point becomes \texttt{int}; a path where any value anywhere
in the corpus is written with a decimal point becomes \texttt{float} for the
whole column.}

\begin{lstlisting}[style=json]
{ "qty": 1 }
{ "qty": 2 }
  -> qty : int

{ "price": 1 }
{ "price": 2.5 }
  -> price : float          (widened: one float anywhere widens the column)

{ "x": 1.0 }
  -> x : float              (token is a float, even though the value is round)
\end{lstlisting}

Widening applies at every path, at any depth. Root fields and array-element
fields obey the same rule; they differ only in how many observations they
receive, since a root field is observed once per document while an
array-element field is observed once per element:

\begin{lstlisting}[style=json]
{ "price": 1,   "items": [ { "qty": 1 } ] }
{ "price": 2.5, "items": [ { "qty": 2 }, { "qty": 3 } ] }
{ "price": 3,   "items": [] }

  .price        sees 1, 2.5, 3   -> float
  .items[].qty  sees 1, 2, 3     -> int
\end{lstlisting}

\rationale{The parser already draws the lexical line --- yojson reports
\texttt{`Int} and \texttt{`Float} as distinct constructors --- so option (b)
costs nothing to implement and respects the distinction the document author
wrote when they wrote it. Widening rather than rejecting keeps inference
monotone: adding documents to a corpus can only loosen a type, never
invalidate one already inferred, which matters because the corpus is not
usually known in full ahead of time.}

\limitation{The lexical signal is not always preserved by the producer.
JavaScript has only doubles, and \texttt{JSON.stringify(1.0)} emits
\texttt{1}, so a corpus of genuinely fractional values that all happen to be
round will infer \texttt{int}. Widening cannot recover this, because the
information is absent from the data. This is a limit of inferring types from
values rather than a consequence of the lexical choice: in any corpus where
some value is fractional, widening gives the right answer.}

\subsection*{S1b. Integers outside native \texttt{int} range}

OCaml's native \texttt{int} is 63 bits on a 64-bit platform. JSON places no
bound on integer literals, so a corpus can carry integers that do not fit.

\begin{lstlisting}[style=json]
{ "id": 12345678901234567890 }
\end{lstlisting}

\options
\begin{enumerate}[label=\textbf{(\alph*)}]
  \item \textbf{Reject}, with an error naming the path and the offending
        value. Keeps the accepted fragment small at the cost of refusing
        corpora other tools accept.
  \item \textbf{Promote the column to \texttt{Int64}.} Covers realistic
        64-bit identifiers exactly; still bounded, so genuinely larger values
        remain an error.
  \item \textbf{Promote the column to \texttt{float}.} \selected
  \item \textbf{Arbitrary precision via zarith} (\texttt{Z.t}). Never lossy,
        but adds a dependency on GMP and makes the generated modules heavier
        to consume.
\end{enumerate}

\decision{Promote to \texttt{float}. A path carrying any integer literal
outside native \texttt{int} range becomes \texttt{float} for the whole
column.}

\begin{lstlisting}[style=json]
{ "id": 1 }
{ "id": 12345678901234567890 }

  -> id : float             (one out-of-range literal widens the column)
\end{lstlisting}

\rationale{This unifies S1b with S1a rather than adding a second mechanism.
There is one widening lattice, \texttt{int} $\sqsubseteq$ \texttt{float}, and
two things push a column up it: a fractional literal, and an integer literal
that does not fit. Inference stays monotone, and no new type appears in the
generated code.}

\implnote{Overflow is detectable at parse time without custom lexing. Yojson's
\texttt{Safe} and \texttt{Raw} parsers report an out-of-range integer literal
as \texttt{`Intlit of string}, carrying the literal rather than failing;
verified against yojson 3.0.0. Parsing with \texttt{Basic}, which lacks
\texttt{`Intlit}, would lose this distinction.}

\limitation{An IEEE double has a 53-bit significand, so integers above
$2^{53}$ are rounded on conversion. A path of 64-bit identifiers will be
silently altered, and two distinct identifiers can collapse to the same
\texttt{float}. This is accepted deliberately; option (b) is the upgrade path
and changes nothing else in this document.}
